\documentclass[11pt]{article}
\usepackage[margin=1in]{geometry}
\usepackage[T1]{fontenc}
\usepackage[utf8]{inputenc}
\usepackage[french]{babel}
\usepackage{amsmath,amssymb,amsthm}
\usepackage[colorlinks=true,linkcolor=blue,urlcolor=blue]{hyperref}
\newtheorem{problem}{Problème}
\newenvironment{refs}{\par\small\noindent\emph{Références :}\begin{itemize}\setlength\itemsep{0pt}}{\end{itemize}\normalsize}
\title{Hors de la Caverne \\ \Large Algèbre abstraite}
\author{}
\date{}
\begin{document}
\maketitle
\noindent\emph{Des problèmes, pas de réponses.} Chaque problème ci-dessous est posé
sans solution. Les références indiquent où trouver les connaissances nécessaires.
\medskip


\section*{Lycée}

\begin{problem}[{Casting out nines and the clock of powers --- Lycée}]
\textbf{ALG-01.} \textbf{Problème.} Two integers $a$ and $b$ are called \textit{congruent
modulo $n$}, written $a \equiv b \pmod{n}$, if $n$ divides $a-b$.
\begin{enumerate}
\item[(a)] Prove that
congruence modulo $n$ respects addition and multiplication: if $a \equiv a' \pmod{n}$ and
$b \equiv b' \pmod{n}$, then $a+b \equiv a'+b' \pmod{n}$ and $ab \equiv a'b' \pmod{n}$.
\item[(b)] Using (i), prove that a positive integer is divisible by 9 if and only if the sum of its
decimal digits is divisible by 9.
\item[(c)] Determine, with proof, the last decimal digit of
$7^{2026}$.
\end{enumerate}

\end{problem}

\noindent Modular arithmetic was made systematic by Gauss in the opening pages of his
\textit{Disquisitiones Arithmeticae} (1801), written when he was 24; the $\equiv$ notation is his. The
digit-sum test ("casting out nines") is far older --- medieval accountants used it to check
ledgers. Behind part (iii) hides the first example of a group in this collection: the residues
modulo 10 that have multiplicative inverses, cycling under multiplication.

\begin{refs}
\item Contexte : Arithmétique modulaire --- Wikipédia (\url{https://fr.wikipedia.org/wiki/Arithm%C3%A9tique_modulaire})
\item Contexte : Disquisitiones arithmeticae --- Wikipédia (\url{https://fr.wikipedia.org/wiki/Disquisitiones_arithmeticae})
\item Lecture : J. Fraleigh, \textit{A First Course in Abstract Algebra}, sections 0 and 4
\end{refs}

\bigskip

\begin{problem}[{The twenty-four rotations of the cube --- Lycée}]
\textbf{ALG-02.} \textbf{Problème.}
\begin{enumerate}
\item[(a)] Prove that a cube has exactly 24 rotational symmetries
--- rotations of space that carry the cube onto itself --- and justify that your list is complete.
\item[(b)] A cube has four space diagonals (segments joining opposite vertices). Show that every
rotational symmetry permutes these four diagonals, and that distinct rotations give distinct
permutations.
\item[(c)] Conclude that the rotations of the cube realize \textit{every} one of the
$4! = 24$ possible rearrangements of the four diagonals.
\end{enumerate}

\end{problem}

\noindent Counting symmetries is where group theory touches the eye. The regular solids
were studied by the Greeks --- Euclid's \textit{Elements} ends with them --- but the insight that their
symmetries form a closed system of composable operations belongs to the nineteenth century, and
part (iii) says the cube's rotation group "is" the symmetric group $S_4$ in disguise.
Klein's 1872 Erlangen program proposed defining geometry itself by its symmetry groups.

\begin{refs}
\item Contexte : Symétrie octaédrique --- Wikipédia (en anglais) (\url{https://en.wikipedia.org/wiki/Octahedral_symmetry})
\item Contexte : Groupe de symétrie --- Wikipédia (\url{https://fr.wikipedia.org/wiki/Groupe_de_sym%C3%A9trie})
\item Lecture : M. Artin, \textit{Algebra}, ch. 6 (Symmetry)
\end{refs}

\bigskip

\begin{problem}[{When is $n^4 + 4$ prime? --- Lycée}]
\textbf{ALG-03.} \textbf{Problème.} Determine, with proof, all positive integers $n$ for which
$n^4 + 4$ is a prime number. More generally, show that $n^4 + 4m^4$ is composite for all
integers $n, m \ge 2$.
\end{problem}

\noindent The key is a factorization that most people insist cannot exist --- a sum
that factors. It is named after Sophie Germain (1776--1831), who taught herself mathematics
from her father's library during the French Revolution and corresponded with Gauss under a
male pseudonym until he learned, with delight, who "Monsieur Le Blanc" really was. The
identity behind this problem appears throughout olympiad mathematics and in her work on
Fermat's Last Theorem.

\begin{refs}
\item Contexte : Identité de Sophie Germain --- Wikipédia (\url{https://fr.wikipedia.org/wiki/Identit%C3%A9_de_Sophie_Germain})
\item Contexte : Sophie Germain --- Wikipédia (\url{https://fr.wikipedia.org/wiki/Sophie_Germain})
\item Lecture : A. Engel, \textit{Problem-Solving Strategies}, ch. 6 (Number Theory)
\end{refs}

\bigskip

\begin{problem}[{The rational root theorem, and why $\sqrt[3]{2}$ is irrational --- Lycée}]
\textbf{ALG-04.} \textbf{Problème.}
\begin{enumerate}
\item[(a)] Prove the rational root theorem: if a polynomial
$a_n x^n + \cdots + a_1 x + a_0$ with integer coefficients has a rational root $p/q$ in
lowest terms, then $p$ divides $a_0$ and $q$ divides $a_n$.
\item[(b)] Deduce that
$\sqrt[3]{2}$ is irrational.
\item[(c)] Show that $x^3 - x - 1$ has no rational root, and
explain why for a cubic this already proves it cannot be factored into polynomials of lower
degree with rational coefficients.
\end{enumerate}

\end{problem}

\noindent This is the humble ancestor of everything in the Graduate band below: the
irrationality of $\sqrt[3]{2}$ is the first, easiest obstruction in the 2000-year-old problem
of doubling the cube (ALG-18). The theorem itself is essentially due to Descartes and was
routine by Euler's time; what changed later is that "has no rational root" grew into the far
deeper notion of \textit{irreducibility}, the atom-concept of field theory.

\begin{refs}
\item Contexte : Théorème de la racine rationnelle --- Wikipédia (\url{https://fr.wikipedia.org/wiki/Racine_%C3%A9vidente})
\item Contexte : Polynôme irréductible --- Wikipédia (\url{https://fr.wikipedia.org/wiki/Polyn%C3%B4me_irr%C3%A9ductible})
\item Lecture : J. Fraleigh, \textit{A First Course in Abstract Algebra}, section 23
\end{refs}

\bigskip

\begin{problem}[{The symmetries of a rectangle --- Lycée, short}]
\textbf{ALG-26.} \textbf{Problème.} List all four symmetries of a non-square rectangle, write out their composition table, and prove that every non-identity element is its own inverse.
\end{problem}

\noindent This is the Klein four-group, the smallest group that is not cyclic, and it can be discovered with a piece of cardboard before the word "group" is ever defined. Building the table by hand and then noticing the pattern is exactly how the abstract axioms were arrived at historically --- the examples came first.

\begin{refs}
\item Contexte : Groupe de Klein --- Wikipédia (\url{https://fr.wikipedia.org/wiki/Groupe_de_Klein})
\end{refs}

\bigskip

\begin{problem}[{A clock that loses numbers --- Lycée, short}]
\textbf{ALG-27.} \textbf{Problème.} In arithmetic modulo $ 12 $, determine exactly which residues have a multiplicative inverse, and prove that your criterion is correct for a general modulus $ n $.
\end{problem}

\noindent Modular arithmetic looks like ordinary arithmetic until division fails, and finding precisely when it fails --- exactly when the residue shares a factor with the modulus --- is the first genuine theorem most people prove about rings. It is also why prime moduli behave so much better, which is the reason cryptography prefers them.

\begin{refs}
\item Contexte : Inverse modulaire --- Wikipédia (\url{https://fr.wikipedia.org/wiki/Inverse_modulaire})
\end{refs}

\bigskip

\begin{problem}[{Parity, and the 14--15 that never comes out --- Lycée}]
\textbf{ALG-34.} \textbf{Problème.} A \textit{transposition} is a permutation that exchanges two elements and fixes
the rest; every permutation of a finite set is a product of transpositions. The \textit{15 puzzle}
is a $4 \times 4$ tray holding tiles $1, 2, \dots, 15$ and one empty square; a move slides a
tile that is horizontally or vertically adjacent to the empty square into it.
\begin{enumerate}
\item[(a)] Prove that if a permutation $\sigma$ is written as a product of transpositions in two
ways, the two numbers of factors have the same parity. So the \textit{sign}
$\operatorname{sgn}(\sigma) = \pm 1$ is well defined; prove that
$\operatorname{sgn}(\sigma\tau) = \operatorname{sgn}(\sigma)\operatorname{sgn}(\tau)$.
\item[(b)] Deduce that for $n \ge 2$ the even permutations form a subgroup of $S_n$ containing
exactly half its elements.
\item[(c)] Read a position of the puzzle as a permutation $\sigma$ of the sixteen squares (the empty
square counting as a sixteenth piece), with the empty square in row $i$ and column $j$.
Prove that
\[ \operatorname{sgn}(\sigma) \cdot (-1)^{\,i+j} \]
is unchanged by every legal move.
\item[(d)] Conclude that the position obtained from the solved one by exchanging the tiles 14 and 15 can
never be reached, and prove that at most half of the $16!$ arrangements are reachable from any
given one.
\end{enumerate}

\end{problem}

\noindent The puzzle was devised around 1874 by Noyes Palmer Chapman, a postmaster in
Canastota, New York, and became a craze in 1880; Sam Loyd, who had nothing to do with inventing
it, later offered $\$1000$ for a solution to the 14--15 position and from 1891 claimed the puzzle
as his own --- a claim demolished by Slocum and Sonneveld's 2006 history. The prize was safe,
because in 1879 William Woolsey Johnson and William E. Story had already published the invariant
of part (c) in the \textit{American Journal of Mathematics}: one of the earliest arguments in which
an algebraic invariant proves that something cannot be done, rather than merely failing to do it.
The sign of a permutation, born in eighteenth-century work on determinants, is the same
homomorphism that produces the alternating group of ALG-11 and, in the end, the unsolvability of
the quintic.

\begin{refs}
\item Contexte : Taquin --- Wikipédia (\url{https://fr.wikipedia.org/wiki/Taquin})
\item Contexte : Signature d'une permutation --- Wikipédia (\url{https://fr.wikipedia.org/wiki/Signature_d%27une_permutation})
\item Article : Wm. W. Johnson \& W. E. Story, ``Notes on the `15' Puzzle'' (1879), American Journal of Mathematics 2, 397--404
\item Article : A. F. Archer, ``A modern treatment of the 15 puzzle'' (1999), American Mathematical Monthly 106, 793--799
\end{refs}

\bigskip

\begin{problem}[{Every group table is a Latin square --- Lycée, short}]
\textbf{ALG-35.} \textbf{Problème.} Prove that in the multiplication table of a finite group,
every element of the group appears exactly once in each row and exactly once in each column.
Then show, by exhibiting one, that a square array with that property need not be the
multiplication table of any group.
\end{problem}

\noindent The table is Cayley's invention: in his 1854 paper on abstract groups he laid
one out and observed exactly this pattern, which is nothing but the cancellation law made visible.
Squares filled so that each symbol occurs once per row and column had already been studied by
Euler in 1782 under the name \textit{Latin squares}, and the gap between the two notions is
instructive --- a Latin square encodes cancellation, but associativity leaves no visible trace in
the grid, which is why the converse fails and why the objects that satisfy only the row-and-column
condition (quasigroups) form a separate subject.

\begin{refs}
\item Contexte : Table de Cayley --- Wikipédia (\url{https://fr.wikipedia.org/wiki/Table_de_Cayley})
\item Contexte : Carré latin --- Wikipédia (\url{https://fr.wikipedia.org/wiki/Carr%C3%A9_latin})
\end{refs}

\bigskip


\section*{Licence}

\begin{problem}[{What is a group, exactly? --- Licence}]
\textbf{ALG-05.} \textbf{Problème.} A \textit{group} is a set $G$ with an operation $*$ that
is associative, has an identity element, and gives every element an inverse. Determine, with
proof, which of the following are groups:
\begin{enumerate}
\item[(a)] the integers under subtraction;
\item[(b)] the positive
rationals under multiplication;
\item[(c)] the set $\{0, 1, \dots, 11\}$ under multiplication mod 12;
\item[(d)] those elements of $\{0, 1, \dots, 11\}$ having a multiplicative inverse mod 12, under
multiplication mod 12;
\item[(e)] invertible $2 \times 2$ real matrices under matrix multiplication.
Then prove from the axioms alone that in any group the identity is unique, each element has a
unique inverse, and $(ab)^{-1} = b^{-1}a^{-1}$.
\end{enumerate}

\end{problem}

\noindent The axioms were distilled slowly, over a century, from three separate springs:
Galois's groups of permutations of roots (1830s), Klein's and Lie's groups of geometric
transformations, and number-theoretic composition laws going back to Gauss. Cayley wrote down an
abstract definition in 1854 and was largely ignored; only around 1900 did the axiomatic viewpoint
win. The exercise of checking which familiar structures satisfy the axioms --- and which fail, and
why --- is how every algebraist first learns what the axioms are actually for.

\begin{refs}
\item Contexte : Groupe (mathématiques) --- Wikipédia (\url{https://fr.wikipedia.org/wiki/Groupe_(math%C3%A9matiques)})
\item Lecture : D. Dummit \& R. Foote, \textit{Abstract Algebra}, 3rd ed., ch. 1
\item Cours : MIT OCW 18.701 Algebra I (\url{https://ocw.mit.edu/courses/18-701-algebra-i-fall-2010/})
\end{refs}

\bigskip

\begin{problem}[{Lagrange's theorem --- Licence}]
\textbf{ALG-06.} \textbf{Problème.} Let $G$ be a finite group and $H$ a subgroup.
\begin{enumerate}
\item[(a)] Prove
Lagrange's theorem: the order (size) of $H$ divides the order of $G$.
\item[(b)] Deduce that the
order of any element of $G$ divides the order of $G$, and that every group of prime order is
cyclic.
\item[(c)] Deduce Fermat's little theorem: if $p$ is prime and $p$ does not divide $a$,
then $a^{p-1} \equiv 1 \pmod{p}$.
\end{enumerate}

\end{problem}

\noindent Lagrange stated a special case in 1770--71 while studying how the value of a
polynomial expression changes when its variables are permuted --- work that predates the definition
of a group by decades. The modern statement and proof, via the beautiful device of cosets slicing
a group into equal-sized pieces, crystallized with Galois and Jordan. It is the first structural
theorem of the subject: a group's size constrains what can live inside it. That Fermat's 1640
observation about primes falls out as a corollary is an early hint of algebra's unifying power.

\begin{refs}
\item Contexte : Théorème de Lagrange (théorie des groupes) --- Wikipédia (\url{https://fr.wikipedia.org/wiki/Th%C3%A9or%C3%A8me_de_Lagrange_sur_les_groupes})
\item Contexte : Petit théorème de Fermat --- Wikipédia (\url{https://fr.wikipedia.org/wiki/Petit_th%C3%A9or%C3%A8me_de_Fermat})
\item Lecture : D. Dummit \& R. Foote, \textit{Abstract Algebra}, 3rd ed., ch. 3
\item Cours : MIT OCW 18.701 Algebra I (\url{https://ocw.mit.edu/courses/18-701-algebra-i-fall-2010/})
\end{refs}

\bigskip

\begin{problem}[{Cayley's theorem: every group is a group of permutations --- Licence}]
\textbf{ALG-07.} \textbf{Problème.}
\begin{enumerate}
\item[(a)] Prove Cayley's theorem: every group $G$ is isomorphic
to a subgroup of the symmetric group $\mathrm{Sym}(G)$ of all permutations of the set $G$.
\item[(b)] The theorem embeds a group of order $n$ into $S_n$, but often one can do much better ---
or not. Show that the quaternion group $Q_8 = \{\pm 1, \pm i, \pm j, \pm k\}$ admits no
faithful action on fewer than 8 points: it is not isomorphic to a subgroup of $S_n$ for any
$n < 8$.
\end{enumerate}

\end{problem}

\noindent Cayley proved (i) in 1854, in the very paper where he first defined abstract
groups --- it was his argument that the abstract definition loses nothing, since every abstract group
is realized by permutations. Part (ii) is the classic counterpoint: the regular representation can
be wildly inefficient, but for $Q_8$ (Hamilton's quaternion units, 1843) it is provably
optimal. The study of minimal faithful permutation degrees is still an active corner of group
theory.

\begin{refs}
\item Contexte : Théorème de Cayley --- Wikipédia (\url{https://fr.wikipedia.org/wiki/Th%C3%A9or%C3%A8me_de_Cayley})
\item Contexte : Groupe des quaternions --- Wikipédia (\url{https://fr.wikipedia.org/wiki/Groupe_des_quaternions})
\item Lecture : J. Rotman, \textit{An Introduction to the Theory of Groups}, 4th ed., ch. 3
\end{refs}

\bigskip

\begin{problem}[{Subgroup lattices: $D_4$ and $\mathbb{Z}/12\mathbb{Z}$ --- Licence}]
\textbf{ALG-08.} \textbf{Problème.} The subgroups of a group, ordered by inclusion, form its
\textit{subgroup lattice}.
\begin{enumerate}
\item[(a)] Find all subgroups of the dihedral group $D_4$ (the 8 symmetries
of a square), prove your list is complete, and draw the lattice. Mark which subgroups are
normal.
\item[(b)] Prove that every subgroup of a cyclic group is cyclic, and that the subgroup
lattice of $\mathbb{Z}/n\mathbb{Z}$ is exactly the lattice of divisors of $n$. Draw it for
$n = 12$.
\item[(c)] Exhibit two nonisomorphic groups with isomorphic subgroup lattices.
\end{enumerate}

\end{problem}

\noindent Drawing a subgroup lattice is to a group theorist what dissection is to an
anatomist. The systematic study of what the lattice does and does not remember about the group was
begun by Ada Rottlaender (a student of Schur) in 1928 and grew into a whole subject, surveyed in
Suzuki's 1956 book; part (iii) shows the lattice is a genuine but lossy invariant. The picture for
$D_4$, with its ten subgroups, reappears in ALG-20 as one half of a Galois correspondence.

\begin{refs}
\item Contexte : Treillis des sous-groupes --- Wikipédia (\url{https://fr.wikipedia.org/wiki/Treillis_des_sous-groupes})
\item Contexte : Groupe diédral --- Wikipédia (\url{https://fr.wikipedia.org/wiki/Groupe_di%C3%A9dral})
\item Lecture : D. Dummit \& R. Foote, \textit{Abstract Algebra}, 3rd ed., section 2.5
\end{refs}

\bigskip

\begin{problem}[{Group actions and Burnside counting --- Licence}]
\textbf{ALG-09.} \textbf{Problème.} A group $G$ \textit{acts} on a set $X$ if each
$g \in G$ permutes $X$ compatibly with the group operation.
\begin{enumerate}
\item[(a)] Prove the orbit--stabilizer
theorem: for finite $G$ and any $x \in X$, $|G| = |\mathrm{Orb}(x)| \cdot |\mathrm{Stab}(x)|$.
\item[(b)] Prove Burnside's lemma: the number of orbits equals
$\frac{1}{|G|} \sum_{g \in G} |\mathrm{Fix}(g)|$, the average number of points fixed by a
group element.
\item[(c)] Use it to count the essentially different ways to paint the six faces of a
cube with three colors, where two paintings are the same if a rotation carries one to the
other.
\end{enumerate}

\end{problem}

\noindent The counting lemma is the subject's most famous misattribution: it appears in
Cauchy (1845) and Frobenius (1887), and Burnside himself cited Frobenius when he included it in his
1897 textbook --- group theorists wryly call it "the lemma that is not Burnside's" (the title of a
1979 note by Peter Neumann). Beyond puzzles, the orbit--stabilizer theorem is the engine of nearly
every counting argument in finite group theory, including the Sylow theorems two problems below.

\begin{refs}
\item Contexte : Lemme de Burnside --- Wikipédia (\url{https://fr.wikipedia.org/wiki/Lemme_de_Burnside})
\item Contexte : Action de groupe --- Wikipédia (\url{https://fr.wikipedia.org/wiki/Action_de_groupe_(math%C3%A9matiques)})
\item Lecture : M. Artin, \textit{Algebra}, ch. 6--7
\item Cours : MIT OCW 18.701 Algebra I (\url{https://ocw.mit.edu/courses/18-701-algebra-i-fall-2010/})
\end{refs}

\bigskip

\begin{problem}[{Sylow theory in action: groups of order $pq$ --- Licence}]
\textbf{ALG-10.} \textbf{Problème.} The Sylow theorems say: if $p^k$ is the largest power of
the prime $p$ dividing $|G|$, then subgroups of order $p^k$ (Sylow $p$-subgroups) exist,
are all conjugate, and their number is $\equiv 1 \pmod{p}$ and divides $|G|$. Take these as
given (or prove them).
\begin{enumerate}
\item[(a)] Prove that every group of order 15 is cyclic.
\item[(b)] More generally,
let $p < q$ be primes with $p$ not dividing $q - 1$; prove every group of order $pq$
is cyclic.
\item[(c)] Show that when $p$ does divide $q - 1$, a nonabelian group of order
$pq$ exists.
\item[(d)] Prove that no group of order 30 is simple.
\end{enumerate}

\end{problem}

\noindent Ludwig Sylow, a Norwegian high-school teacher for most of his career, published
his three theorems in 1872 in a ten-page paper that became the cornerstone of finite group theory.
They are the partial converse Lagrange's theorem lacks: divisibility of the order forces existence
of subgroups. Arguments like (i)--(iv) --- juggling counts of Sylow subgroups until a group is forced
to be cyclic, or forced to have a normal subgroup --- are the daily bread of the classification
programme that culminates in ALG-22.

\begin{refs}
\item Contexte : Théorèmes de Sylow --- Wikipédia (\url{https://fr.wikipedia.org/wiki/Th%C3%A9or%C3%A8mes_de_Sylow})
\item Lecture : D. Dummit \& R. Foote, \textit{Abstract Algebra}, 3rd ed., ch. 4
\item Cours : MIT OCW 18.701 Algebra I (\url{https://ocw.mit.edu/courses/18-701-algebra-i-fall-2010/})
\end{refs}

\bigskip

\begin{problem}[{The simplicity of $A_5$ --- Licence}]
\textbf{ALG-11.} \textbf{Problème.} A group is \textit{simple} if its only normal subgroups are
itself and the trivial group.
\begin{enumerate}
\item[(a)] Prove that the alternating group $A_5$ (even permutations
of five letters, order 60) is simple.
\item[(b)] Prove that every nonabelian simple group of order at
most 60 has order exactly 60, so $A_5$ is the smallest nonabelian simple group.
\item[(c)] Prove
that $A_5$ is isomorphic to the rotation group of the icosahedron.
\end{enumerate}

\end{problem}

\noindent This one group is the villain --- or hero --- of the whole story. Galois knew by
1832 that $A_5$ is simple, and its simplicity is precisely what makes the general quintic
unsolvable by radicals (ALG-21). Part (ii), a classic tour of Sylow counting, explains why degree
five is the first bad degree. Part (iii), beloved of Klein (whose 1884 \textit{Lectures on the
Icosahedron} built a whole theory around it), says the obstruction has been sitting inside a
Platonic solid all along.

\begin{refs}
\item Contexte : Groupe alterné --- Wikipédia (\url{https://fr.wikipedia.org/wiki/Groupe_altern%C3%A9})
\item Contexte : Groupe simple --- Wikipédia (\url{https://fr.wikipedia.org/wiki/Groupe_simple})
\item Lecture : D. Dummit \& R. Foote, \textit{Abstract Algebra}, 3rd ed., sections 4.5--4.6
\end{refs}

\bigskip

\begin{problem}[{Semidirect products --- Licence}]
\textbf{ALG-12.} \textbf{Problème.} Given groups $N$ and $H$ and a homomorphism
$\varphi : H \to \operatorname{Aut}(N)$, the \textit{semidirect product}
$N \rtimes_\varphi H$ is the set $N \times H$ with multiplication
$(n_1, h_1)(n_2, h_2) = (n_1 \, \varphi(h_1)(n_2), \; h_1 h_2)$.
\begin{enumerate}
\item[(a)] Verify this is a group
containing $N$ as a normal subgroup.
\item[(b)] Show that the dihedral group $D_n$ is a
semidirect product $\mathbb{Z}/n\mathbb{Z} \rtimes \mathbb{Z}/2\mathbb{Z}$.
\item[(c)] Construct
a nonabelian group of order 21, and prove that up to isomorphism there is exactly one.
\item[(d)] For primes $p < q$, determine exactly when a nonabelian semidirect product
$\mathbb{Z}/q\mathbb{Z} \rtimes \mathbb{Z}/p\mathbb{Z}$ exists.
\end{enumerate}

\end{problem}

\noindent The semidirect product is the algebraist's basic tool for building new groups
from old, formalized by Hölder in the 1890s during the first classification of groups of small
order. It answers the "extension problem" in its easiest case: how can a group be assembled from a
normal subgroup and a quotient? Hölder's program --- understand simple groups, then understand how
they glue --- is still the blueprint of finite group theory.

\begin{refs}
\item Contexte : Produit semi-direct --- Wikipédia (\url{https://fr.wikipedia.org/wiki/Produit_semi-direct})
\item Lecture : D. Dummit \& R. Foote, \textit{Abstract Algebra}, 3rd ed., section 5.5
\item Lecture : J. Rotman, \textit{An Introduction to the Theory of Groups}, 4th ed., ch. 7
\end{refs}

\bigskip

\begin{problem}[{Rings, ideals, and what quotients detect --- Licence}]
\textbf{ALG-13.} \textbf{Problème.} Let $R$ be a commutative ring with 1. An \textit{ideal}
$I \subset R$ is a nonempty subset closed under addition and under multiplication by arbitrary
ring elements.
\begin{enumerate}
\item[(a)] Prove that the ideals of $R$ are exactly the kernels of ring homomorphisms
out of $R$.
\item[(b)] Prove that the quotient $R/I$ is a field if and only if $I$ is maximal
(contained in no larger proper ideal), and an integral domain if and only if $I$ is prime
($ab \in I$ forces $a \in I$ or $b \in I$).
\item[(c)] Determine all prime and all maximal
ideals of $\mathbb{Z}$, and of the polynomial ring $\mathbb{R}[x]$.
\end{enumerate}

\end{problem}

\noindent Ideals were born of a crisis: Kummer's "ideal numbers" (1840s) rescued unique
factorization in the number rings where it fails, and Dedekind recast them (1871) as the subsets
now bearing the name. Emmy Noether's 1921 paper \textit{Idealtheorie in Ringbereichen} made ideals
the central objects of commutative algebra --- the modern subject is unrecognizable without her.
The dictionary in (ii), translating properties of quotients into properties of ideals, is the
template for algebraic geometry's dictionary between rings and spaces.

\begin{refs}
\item Contexte : Idéal (théorie des anneaux) --- Wikipédia (\url{https://fr.wikipedia.org/wiki/Id%C3%A9al})
\item Contexte : Anneau (mathématiques) --- Wikipédia (\url{https://fr.wikipedia.org/wiki/Anneau_(math%C3%A9matiques)})
\item Lecture : D. Dummit \& R. Foote, \textit{Abstract Algebra}, 3rd ed., ch. 7
\item Cours : MIT OCW 18.702 Algebra II (\url{https://ocw.mit.edu/courses/18-702-algebra-ii-spring-2011/})
\end{refs}

\bigskip

\begin{problem}[{Quotient constructions: building $\mathbb{C}$ from $\mathbb{R}[x]$ --- Licence}]
\textbf{ALG-14.} \textbf{Problème.}
\begin{enumerate}
\item[(a)] Prove the first isomorphism theorem, for groups and for
rings: if $f : A \to B$ is a homomorphism, then $A/\ker f \cong \operatorname{im} f$.
\item[(b)] Prove that $\mathbb{R}[x]/(x^2 + 1)$ is a field isomorphic to the complex numbers
$\mathbb{C}$.
\item[(c)] More generally, prove that if $p(x)$ is irreducible over a field $F$,
then $F[x]/(p(x))$ is a field containing (a copy of) $F$ in which $p$ has a root.
\item[(d)] Use (iii) to build a field with exactly 4 elements, and write out its addition and
multiplication tables.
\end{enumerate}

\end{problem}

\noindent Part (iii) is Kronecker's construction (1887): given any polynomial, a root can
be \textit{manufactured} by pure algebra, with no appeal to $\mathbb{C}$ or to analysis. This
inverted the history --- mathematicians had adjoined $i$ to the reals with a guilty conscience for
three centuries, and here was a machine that mints roots legally. It is the engine that makes
field theory self-sufficient and the direct route to the finite fields of ALG-17.

\begin{refs}
\item Contexte : Anneau quotient --- Wikipédia (\url{https://fr.wikipedia.org/wiki/Anneau_quotient})
\item Contexte : Théorèmes d'isomorphisme --- Wikipédia (\url{https://fr.wikipedia.org/wiki/Th%C3%A9or%C3%A8mes_d%27isomorphisme})
\item Lecture : M. Artin, \textit{Algebra}, ch. 11
\item Cours : MIT OCW 18.702 Algebra II (\url{https://ocw.mit.edu/courses/18-702-algebra-ii-spring-2011/})
\end{refs}

\bigskip

\begin{problem}[{Eisenstein's criterion and the cyclotomic polynomial --- Licence}]
\textbf{ALG-15.} \textbf{Problème.}
\begin{enumerate}
\item[(a)] Prove Gauss's lemma: a polynomial with integer
coefficients that factors over $\mathbb{Q}$ already factors over $\mathbb{Z}$ (into factors
of the same degrees).
\item[(b)] Prove Eisenstein's criterion: if a prime $p$ divides every
coefficient of $f(x) = a_n x^n + \cdots + a_0$ except $a_n$, and $p^2$ does not divide
$a_0$, then $f$ is irreducible over $\mathbb{Q}$.
\item[(c)] Deduce that $x^n - p$ is
irreducible over $\mathbb{Q}$ for every $n \ge 1$ and prime $p$.
\item[(d)] Show that for
$p$ prime, $\Phi_p(x) = x^{p-1} + x^{p-2} + \cdots + 1$ is irreducible over $\mathbb{Q}$.
(The standard trick: apply the criterion after a change of variable.)
\end{enumerate}

\end{problem}

\noindent Irreducibility is easy to define and maddening to verify; Eisenstein's 1850
criterion is the one tool every algebraist carries. It was actually published four years earlier by
Schönemann, and historians now write "Schönemann--Eisenstein"; both were after part (iv), the
irreducibility of the cyclotomic polynomial, which Gauss had needed for his 17-gon (ALG-19).
Eisenstein died at 29; Gauss reportedly ranked him with Archimedes and Newton.

\begin{refs}
\item Contexte : Critère d'Eisenstein --- Wikipédia (\url{https://fr.wikipedia.org/wiki/Crit%C3%A8re_d%27Eisenstein})
\item Contexte : Lemme de Gauss (polynômes) --- Wikipédia (\url{https://fr.wikipedia.org/wiki/Lemme_de_Gauss_(polyn%C3%B4mes)})
\item Lecture : D. Dummit \& R. Foote, \textit{Abstract Algebra}, 3rd ed., ch. 9
\end{refs}

\bigskip

\begin{problem}[{Field extensions and the tower law --- Licence}]
\textbf{ALG-16.} \textbf{Problème.} If $F \subset K$ are fields, $K$ is a vector space over
$F$; its dimension is the \textit{degree} $[K : F]$.
\begin{enumerate}
\item[(a)] Prove the tower law: if
$F \subset K \subset L$ then $[L : F] = [L : K] \cdot [K : F]$.
\item[(b)] Prove that
$[\mathbb{Q}(\sqrt[3]{2}) : \mathbb{Q}] = 3$ and
$[\mathbb{Q}(\sqrt{2}, \sqrt{3}) : \mathbb{Q}] = 4$, and find a single element $\alpha$
with $\mathbb{Q}(\sqrt{2}, \sqrt{3}) = \mathbb{Q}(\alpha)$.
\item[(c)] Prove that the elements of
an extension of $F$ that are algebraic over $F$ (roots of polynomials with coefficients in
$F$) form a field. Conclude that the sum of two algebraic numbers is algebraic --- without ever
exhibiting its polynomial.
\end{enumerate}

\end{problem}

\noindent Measuring a field extension by a whole number is one of the great compressions
in mathematics, and it is linear algebra smuggled into arithmetic: dimension arguments do all the
work. The viewpoint matured from Dedekind through Steinitz, whose 1910 paper \textit{Algebraische
Theorie der Körper} organized all of field theory; the tower law's multiplicativity is exactly
what powers the impossibility proofs of ALG-18.

\begin{refs}
\item Contexte : Extension de corps --- Wikipédia (\url{https://fr.wikipedia.org/wiki/Extension_de_corps})
\item Contexte : Degré d'une extension de corps --- Wikipédia (\url{https://fr.wikipedia.org/wiki/Degr%C3%A9_d%27une_extension})
\item Lecture : I. Stewart, \textit{Galois Theory}, 4th ed., ch. 4--6
\item Cours : MIT OCW 18.702 Algebra II (\url{https://ocw.mit.edu/courses/18-702-algebra-ii-spring-2011/})
\end{refs}

\bigskip

\begin{problem}[{Finite fields --- Licence}]
\textbf{ALG-17.} \textbf{Problème.}
\begin{enumerate}
\item[(a)] Prove that every finite field has $p^n$ elements for
some prime $p$ and $n \ge 1$.
\item[(b)] Prove existence and uniqueness: for each such $p^n$
there is a field of that order, unique up to isomorphism --- the field $\mathbb{F}_{p^n}$,
obtained as the splitting field of $x^{p^n} - x$ over $\mathbb{F}_p$.
\item[(c)] Prove that the
multiplicative group of any finite field is cyclic.
\item[(d)] Show that $\mathbb{F}_{p^m}$ embeds
in $\mathbb{F}_{p^n}$ if and only if $m$ divides $n$, so the subfields of
$\mathbb{F}_{p^n}$ form a copy of the divisor lattice of $n$.
\end{enumerate}

\end{problem}

\noindent Galois introduced these fields in 1830, three years after Abel's death and two
before his own, to study equations modulo $p$ --- they are still called Galois fields, and
E. H. Moore proved the uniqueness (ii) in 1893. A structure invented in the purest corner of
algebra now runs the physical world: every QR code, CD, and deep-space transmission is protected
by Reed--Solomon codes doing arithmetic in $\mathbb{F}_{256}$, and elliptic-curve cryptography
lives over large prime fields.

\begin{refs}
\item Contexte : Corps fini --- Wikipédia (\url{https://fr.wikipedia.org/wiki/Corps_fini})
\item Lecture : D. Dummit \& R. Foote, \textit{Abstract Algebra}, 3rd ed., section 14.3
\item Lecture : R. Lidl \& H. Niederreiter, \textit{Finite Fields}, ch. 2
\end{refs}

\bigskip

\begin{problem}[{A group with no proper subgroup --- Licence, short}]
\textbf{ALG-28.} \textbf{Problème.} Determine all finite groups whose only subgroups are the trivial group and the whole group, and prove your classification is complete.
\end{problem}

\noindent The answer is short --- the cyclic groups of prime order --- and the proof is a clean application of Lagrange's theorem together with the fact that every element generates a subgroup. It is a good first taste of a classification result: not "here is an example" but "here is everything there is".

\begin{refs}
\item Contexte : Théorème de Lagrange --- Wikipédia (\url{https://fr.wikipedia.org/wiki/Th%C3%A9or%C3%A8me_de_Lagrange_sur_les_groupes})
\end{refs}

\bigskip

\begin{problem}[{Centre of a $ p $-group --- Licence, short}]
\textbf{ALG-29.} \textbf{Problème.} Let $ G $ be a finite group of order $ p^n $ with $ p $ prime and $ n \geq 1 $. Prove that the centre of $ G $ is non-trivial, and deduce that every group of order $ p^2 $ is abelian.
\end{problem}

\noindent The proof runs through the class equation, and it is the moment when counting conjugacy classes stops being bookkeeping and starts producing structure out of nothing but the order of the group. The corollary is striking: an arithmetic fact about $ |G| $ forces a commutativity fact about its elements.

\begin{refs}
\item Contexte : Centre d'un groupe --- Wikipédia (\url{https://fr.wikipedia.org/wiki/Centre_d%27un_groupe})
\item Contexte : Équation aux classes --- Wikipédia (\url{https://fr.wikipedia.org/wiki/Action_par_conjugaison})
\end{refs}

\bigskip

\begin{problem}[{Every finitely generated abelian group, listed once --- Licence}]
\textbf{ALG-36.} \textbf{Problème.} Let $G$ be an abelian group generated by finitely many elements.
\begin{enumerate}
\item[(a)] Prove the invariant factor form: there are an integer $r \ge 0$ and integers
$d_1 \mid d_2 \mid \cdots \mid d_k$ with each $d_i \ge 2$ such that
\[ G \;\cong\; \mathbb{Z}^{r} \oplus \mathbb{Z}/d_1\mathbb{Z} \oplus \cdots \oplus \mathbb{Z}/d_k\mathbb{Z}, \]
and prove that $r$ and the list $d_1, \dots, d_k$ are determined by $G$ up to isomorphism.
\item[(b)] Prove the equivalence with the elementary divisor form, in which every summand is cyclic of
prime power order, and explain precisely where the Chinese remainder theorem does the
translating.
\item[(c)] Determine, with proof, all abelian groups of order $360$ up to isomorphism, and prove that
no two groups on your list are isomorphic.
\item[(d)] Prove that for a prime $p$, the number of abelian groups of order $p^n$ up to isomorphism
equals the number of partitions of $n$.
\end{enumerate}

\end{problem}

\noindent This theorem is older than the word "group" in its modern sense. Gauss handled a
first case in the \textit{Disquisitiones} (1801), in the guise of classes of binary quadratic forms;
Kronecker proved it for finite abelian groups in 1870; H. J. S. Smith's 1861 normal form for
integer matrices gives the finitely presented case, since an integer matrix is exactly a
presentation; and Emmy Noether, in 1926, put it in the form used today. Its real weight is that it
is a \textit{complete} classification --- the rarest kind of theorem --- and that the same argument, run
over a polynomial ring instead of $\mathbb{Z}$, produces the rational and Jordan canonical forms
of linear algebra (LIN-14). Uniqueness in part (a) is the subtle half; the existence half is
essentially Gaussian elimination over $\mathbb{Z}$.

\begin{refs}
\item Contexte : Groupe abélien de type fini --- Wikipédia (\url{https://fr.wikipedia.org/wiki/Groupe_ab%C3%A9lien_de_type_fini})
\item Contexte : Théorème des facteurs invariants --- Wikipédia (\url{https://fr.wikipedia.org/wiki/Th%C3%A9or%C3%A8me_des_facteurs_invariants})
\item Lecture : D. Dummit \& R. Foote, \textit{Abstract Algebra}, 3rd ed., section 5.2 and ch. 12
\item Lecture : M. Artin, \textit{Algebra}, ch. 14
\end{refs}

\bigskip

\begin{problem}[{Euclidean, principal, unique --- and the gaps between --- Licence}]
\textbf{ALG-37.} \textbf{Problème.} Let $R$ be an integral domain. Call $R$ \textit{Euclidean} if there is a
function $f : R \setminus \{0\} \to \mathbb{Z}_{\ge 0}$ such that any $a, b$ with $b \ne 0$
admit $a = qb + r$ with $r = 0$ or $f(r) < f(b)$; a \textit{PID} if every ideal is generated
by a single element; a \textit{UFD} if every nonzero nonunit factors into irreducibles, uniquely up
to order and unit multiples.
\begin{enumerate}
\item[(a)] Prove that Euclidean $\Rightarrow$ PID $\Rightarrow$ UFD. Prove also that in a PID every
irreducible element is prime, and that a PID has no infinite strictly increasing chain of
ideals.
\item[(b)] Prove that $\mathbb{Z}[x]$ is a UFD but not a PID, by exhibiting an ideal that no single
element generates.
\item[(c)] Prove that $\mathbb{Z}[\sqrt{-5}]$ is not a UFD, using the norm
$N(a + b\sqrt{-5}) = a^2 + 5b^2$; then say exactly which implication of part (a) this does
\textit{not} contradict, and why.
\item[(d)] Call $u \in R$ a \textit{universal side divisor} if $u$ is a nonzero nonunit and every
$x \in R$ satisfies $u \mid x - r$ for some $r$ that is zero or a unit. Prove that a
Euclidean domain which is not a field must contain one, and prove that the ring
$\mathbb{Z}\!\left[\tfrac{1 + \sqrt{-19}}{2}\right]$ contains none --- so it is not Euclidean,
although it is known to be a PID.
\end{enumerate}

\end{problem}

\noindent Every schoolchild's division-with-remainder is doing more work than it looks:
it is the reason $\mathbb{Z}$ has unique factorization at all. When Kummer found in the 1840s
that unique factorization fails in the cyclotomic rings he needed for Fermat's Last Theorem, the
whole hierarchy above had to be untangled, and Dedekind's ideals (ALG-13) were the repair.
Part (d) is the delicate end of the story: "Euclidean" asserts that some measuring function
exists, so proving a ring is \textit{not} Euclidean requires an obstruction that no function can
dodge. Motzkin supplied one in 1949 with the universal side divisor, and
$\mathbb{Z}\!\left[\tfrac{1+\sqrt{-19}}{2}\right]$ became the first known principal ideal
domain in which no Euclidean algorithm can be installed.

\begin{refs}
\item Contexte : Anneau euclidien --- Wikipédia (\url{https://fr.wikipedia.org/wiki/Anneau_euclidien})
\item Contexte : Anneau principal --- Wikipédia (\url{https://fr.wikipedia.org/wiki/Anneau_principal})
\item Lecture : D. Dummit \& R. Foote, \textit{Abstract Algebra}, 3rd ed., ch. 8
\item Article : Th. Motzkin, ``The Euclidean algorithm'' (1949), Bulletin of the American Mathematical Society 55, 1142--1146
\end{refs}

\bigskip

\begin{problem}[{The largest abelian quotient --- Licence, short}]
\textbf{ALG-38.} \textbf{Problème.} For a group $G$, let $[G,G]$ be the subgroup generated by
all commutators $xyx^{-1}y^{-1}$; prove that $[G,G]$ is normal, that $G/[G,G]$ is abelian,
and that every homomorphism from $G$ to an abelian group factors uniquely through the quotient
map $G \to G/[G,G]$. Determine $[S_n, S_n]$, with proof.
\end{problem}

\noindent The quotient $G/[G,G]$, the \textit{abelianisation}, is the largest piece of a
group that commutative arithmetic can see, and the factoring property says it is the universal
such piece --- a first, painless example of a universal property. Iterating the construction gives
the derived series, whose termination is exactly solvability (ALG-21), so this two-line definition
is what stands between the quartic and the quintic. In topology the same group reappears as the
first homology of a space, by Poincaré's and Hurewicz's theorem that $H_1$ is the abelianisation
of the fundamental group.

\begin{refs}
\item Contexte : Groupe dérivé --- Wikipédia (\url{https://fr.wikipedia.org/wiki/Groupe_d%C3%A9riv%C3%A9})
\item Lecture : D. Dummit \& R. Foote, \textit{Abstract Algebra}, 3rd ed., ch. 6
\end{refs}

\bigskip


\section*{Master}

\begin{problem}[{Doubling the cube, trisecting the angle --- Master}]
\textbf{ALG-18.} \textbf{Problème.} Call a real number \textit{constructible} if it is a
coordinate of a point obtainable from $(0,0)$ and $(1,0)$ by finitely many
straightedge-and-compass steps.
\begin{enumerate}
\item[(a)] Prove that the constructible numbers form a field closed
under square roots, and that every constructible number lies in a tower of fields
$\mathbb{Q} = F_0 \subset F_1 \subset \cdots \subset F_k$ with each
$[F_{i+1} : F_i] = 2$.
\item[(b)] Conclude that every constructible number has degree $2^m$ over
$\mathbb{Q}$ for some $m$.
\item[(c)] Prove that $\sqrt[3]{2}$ is not constructible: the cube
cannot be doubled.
\item[(d)] Prove that $\cos 20^\circ$ has degree 3 over $\mathbb{Q}$, and
conclude that the $60^\circ$ angle --- hence a general angle --- cannot be trisected.
\end{enumerate}

\end{problem}

\noindent The Delian problem (legend: the oracle at Delos demanded an altar of double the
volume) and the trisection tormented geometers for 2200 years. The impossibility proofs came only
in 1837, from Pierre Wantzel, a 23-year-old engineer, in a seven-page paper in Liouville's journal
--- and were so far ahead of their audience that Wantzel remained obscure for a century. The proof
is a two-line consequence of the tower law: powers of two cannot make a three. It is the purest
example of algebra answering geometry's oldest questions.

\begin{refs}
\item Contexte : Duplication du cube --- Wikipédia (\url{https://fr.wikipedia.org/wiki/Duplication_du_cube})
\item Contexte : Trisection de l'angle --- Wikipédia (\url{https://fr.wikipedia.org/wiki/Trisection_de_l%27angle})
\item Lecture : I. Stewart, \textit{Galois Theory}, 4th ed., ch. 7
\item Article : P. Wantzel, ``Recherches sur les moyens de reconnaître si un problème de Géométrie peut se résoudre avec la règle et le compas'' (1837), J. Math. Pures Appl. 2, 366--372
\end{refs}

\bigskip

\begin{problem}[{Gauss's 17-gon and constructible polygons --- Master}]
\textbf{ALG-19.} \textbf{Problème.}
\begin{enumerate}
\item[(a)] Show that the regular $n$-gon is constructible if and
only if $\cos(2\pi/n)$ is a constructible number.
\item[(b)] Prove that the regular 17-gon is
constructible.
\item[(c)] Prove that the regular 7-gon and 9-gon are not.
\item[(d)] Prove the
Gauss--Wantzel theorem: the regular $n$-gon is constructible if and only if $n$ is a power of
2 times a product of distinct Fermat primes (primes of the form $2^{2^k} + 1$). You may use
the irreducibility of cyclotomic polynomials.
\end{enumerate}

\end{problem}

\noindent On 30 March 1796, the eighteen-year-old Gauss woke to see how to construct the
17-gon --- the first advance on ruler-and-compass polygons since antiquity --- and the discovery decided
him for mathematics over philology. He kept a mathematical diary from that day on, and asked for a
17-gon on his tombstone (the stonemason refused: it would look like a circle). Sufficiency is
Gauss's (Disquisitiones, \S{}VII); necessity was proved by Wantzel in 1837. Whether infinitely many
Fermat primes exist --- only 3, 5, 17, 257, 65537 are known --- remains open, so nobody knows whether
the list of constructible polygons with an odd number of sides is finite.

\begin{refs}
\item Contexte : Polygone constructible --- Wikipédia (\url{https://fr.wikipedia.org/wiki/Polygone_constructible})
\item Contexte : Heptadécagone --- Wikipédia (\url{https://fr.wikipedia.org/wiki/Heptad%C3%A9cagone})
\item Contexte : Nombre de Fermat --- Wikipédia (\url{https://fr.wikipedia.org/wiki/Nombre_de_Fermat})
\item Lecture : M. Artin, \textit{Algebra}, ch. 15--16
\end{refs}

\bigskip

\begin{problem}[{The Galois correspondence --- Master}]
\textbf{ALG-20.} \textbf{Problème.} Let $K/F$ be a finite Galois extension (normal and
separable) with Galois group $G = \mathrm{Gal}(K/F)$, the group of automorphisms of $K$
fixing $F$.
\begin{enumerate}
\item[(a)] Prove the fundamental theorem of Galois theory: the maps $H \mapsto K^H$
(fixed field) and $L \mapsto \mathrm{Gal}(K/L)$ are mutually inverse, inclusion-reversing
bijections between subgroups of $G$ and intermediate fields $F \subset L \subset K$;
moreover $[K : K^H] = |H|$, and $H$ is normal in $G$ exactly when $K^H/F$ is Galois,
with $\mathrm{Gal}(K^H/F) \cong G/H$.
\item[(b)] Work the correspondence out completely for the
splitting field of $x^3 - 2$ over $\mathbb{Q}$: find the Galois group, all subgroups, and
all intermediate fields.
\item[(c)] Do the same for $\mathbb{Q}(\zeta_{17})$, the field of 17th
roots of unity, and explain what your lattice has to do with ALG-19.
\end{enumerate}

\end{problem}

\noindent This is the theorem Galois sketched in the memoir he left the night before his
fatal duel in May 1832, aged 20 --- rejected earlier by the Academy, deciphered by Liouville in 1846,
and given its modern field-theoretic form by Dedekind, Weber, and Artin. It is mathematics' first
great \textit{duality}: a dictionary under which hard questions about fields become finite questions
about groups. Everything after it --- from the unsolvability of the quintic to the Langlands program
--- speaks its language.

\begin{refs}
\item Contexte : Théorème fondamental de la théorie de Galois --- Wikipédia (\url{https://fr.wikipedia.org/wiki/Th%C3%A9or%C3%A8me_fondamental_de_la_th%C3%A9orie_de_Galois})
\item Contexte : Théorie de Galois --- Wikipédia (\url{https://fr.wikipedia.org/wiki/Th%C3%A9orie_de_Galois})
\item Lecture : I. Stewart, \textit{Galois Theory}, 4th ed., ch. 8--12
\item Cours : MIT OCW 18.702 Algebra II (\url{https://ocw.mit.edu/courses/18-702-algebra-ii-spring-2011/})
\end{refs}

\bigskip

\begin{problem}[{The insolvability of the general quintic --- Master}]
\textbf{ALG-21.} \textbf{Problème.} A group $G$ is \textit{solvable} if it has a chain of
subgroups from $G$ down to the trivial group, each normal in the next, with abelian quotients.
\begin{enumerate}
\item[(a)] Prove that $S_n$ is solvable for $n \le 4$ and not solvable for $n \ge 5$ (use
ALG-11).
\item[(b)] Prove Galois's criterion, in characteristic 0: if a polynomial is solvable by
radicals --- its roots expressible from the coefficients by field operations and $k$-th roots ---
then its Galois group is solvable.
\item[(c)] Prove that the Galois group of
$f(x) = x^5 - 4x + 2$ over $\mathbb{Q}$ is all of $S_5$. (Show $f$ is irreducible with
exactly three real roots, and use what complex conjugation and an element of order 5 must be.)
(iv) Conclude that $f = 0$ --- and hence the general quintic --- cannot be solved by radicals.
\end{enumerate}

\end{problem}

\noindent Three centuries separate the cubic and quartic formulas of dal Ferro, Tartaglia,
Cardano, and Ferrari (1500s) from the proof that degree five is different. Ruffini published a
flawed impossibility proof in 1799; Abel, in poverty, closed the gap in 1824 (he died at 26);
Galois explained \textit{why}, attaching a group to each equation and reading solvability off the
group. The result is often the first genuine impossibility theorem a mathematician meets: not
"no one has found a formula," but "the symmetry group forbids one."

\begin{refs}
\item Contexte : Théorème d'Abel--Ruffini --- Wikipédia (\url{https://fr.wikipedia.org/wiki/Th%C3%A9or%C3%A8me_d%27Abel_(alg%C3%A8bre)})
\item Contexte : Groupe résoluble --- Wikipédia (\url{https://fr.wikipedia.org/wiki/Groupe_r%C3%A9soluble})
\item Lecture : I. Stewart, \textit{Galois Theory}, 4th ed., ch. 13--15
\item Lecture : D. Dummit \& R. Foote, \textit{Abstract Algebra}, 3rd ed., section 14.7
\end{refs}

\bigskip

\begin{problem}[{Reading the classification of finite simple groups --- Master}]
\textbf{ALG-22.} \textbf{Problème.} The classification of finite simple groups (CFSG) states:
every finite simple group is cyclic of prime order, an alternating group $A_n$ ($n \ge 5$),
a group of Lie type (16 families of matrix-like groups over finite fields), or one of exactly
26 sporadic groups, the largest being the Monster, of order about $8 \times 10^{53}$. This is
a reading-and-verification problem.
\begin{enumerate}
\item[(a)] Prove the easy fragment: every abelian finite simple
group is $\mathbb{Z}/p\mathbb{Z}$.
\item[(b)] Prove that simple groups are the "atoms" of finite
group theory, by proving the Jordan--Hölder theorem: any two composition series of a finite
group have the same simple factors.
\item[(c)] Locate in the classification: $A_5$, the group
$\mathrm{PSL}_2(\mathbb{F}_7)$ of order 168, and the Mathieu group $M_{23}$.
\item[(d)] Write a
two-page account, with citations, of what was proved when: the odd-order theorem
(Feit--Thompson, 1963), the announcement of completion (1983), the quasithin gap closed by
Aschbacher and Smith (2004), and the ongoing second-generation proof of Gorenstein, Lyons, and
Solomon.
\end{enumerate}

\end{problem}

\noindent The CFSG is the largest proof in mathematical history: tens of thousands of
journal pages by about a hundred authors over five decades. No single person has read all of it,
which raises honest philosophical questions about mathematical certainty that the second-generation
rewrite (planned at roughly a dozen volumes) is meant to settle. A working algebraist today mostly
uses CFSG as a black box --- which makes knowing exactly what it says, and what it cost, part of
basic literacy.

\begin{refs}
\item Contexte : Classification des groupes finis simples --- Wikipédia (\url{https://fr.wikipedia.org/wiki/Classification_des_groupes_finis_simples})
\item Contexte : Liste des groupes finis simples --- Wikipédia (\url{https://fr.wikipedia.org/wiki/Liste_des_groupes_finis_simples})
\item Lecture : R. Wilson, \textit{The Finite Simple Groups} (Springer GTM 251)
\item Article : R. Solomon, ``A brief history of the classification of the finite simple groups'' (2001), Bull. Amer. Math. Soc. 38, 315--352
\end{refs}

\bigskip

\begin{problem}[{Irreducible by Eisenstein, after a shift --- Master, short}]
\textbf{ALG-30.} \textbf{Problème.} Prove that the cyclotomic polynomial $ \Phi_p(x) = x^{p-1} + \cdots + x + 1 $ is irreducible over $ \mathbb{Q} $ for prime $ p $, by applying Eisenstein's criterion to $ \Phi_p(x+1) $.
\end{problem}

\noindent Eisenstein's criterion does not apply to $ \Phi_p $ directly, and the substitution $ x \mapsto x+1 $ --- which changes nothing about irreducibility --- makes it apply immediately. Learning that a change of variable can turn an inapplicable theorem into an applicable one is a technique you will use for the rest of your mathematical life.

\begin{refs}
\item Contexte : Polynôme cyclotomique --- Wikipédia (\url{https://fr.wikipedia.org/wiki/Polyn%C3%B4me_cyclotomique})
\item Contexte : Critère d'Eisenstein --- Wikipédia (\url{https://fr.wikipedia.org/wiki/Crit%C3%A8re_d%27Eisenstein})
\end{refs}

\bigskip

\begin{problem}[{The character table of $ S_3 $ --- Master, short}]
\textbf{ALG-31.} \textbf{Problème.} Construct the character table of the symmetric group $ S_3 $ from scratch: identify the three conjugacy classes, find the three irreducible complex representations, and verify the orthogonality relations for the rows.
\end{problem}

\noindent This is the smallest non-abelian character table and the standard entry point to representation theory, where a group is studied by the matrices it can act as. The orthogonality check is the payoff: the rows behave like an orthonormal basis, which is why characters determine a representation completely.

\begin{refs}
\item Contexte : Table de caractères --- Wikipédia (\url{https://fr.wikipedia.org/wiki/Table_de_caract%C3%A8res_(math%C3%A9matiques)})
\item Lecture : Jean-Pierre Serre, \textit{Linear Representations of Finite Groups}, ch. 2
\end{refs}

\bigskip

\begin{problem}[{Free groups, presentations, and where the subgroups are --- Master}]
\textbf{ALG-39.} \textbf{Problème.} Let $S$ be a set. The \textit{free group} $F(S)$ is the group of reduced
words in the letters of $S$ and their formal inverses, multiplied by concatenation followed by
cancellation.
\begin{enumerate}
\item[(a)] Prove that $F(S)$ has the universal property: any map from $S$ into a group $H$ extends
to a unique homomorphism $F(S) \to H$. Deduce that every group is a quotient of a free group, so
every group has a presentation $\langle S \mid R \rangle$.
\item[(b)] Prove von Dyck's theorem in a concrete case: show that
$\langle r, s \mid r^{n}, s^{2}, (rs)^{2} \rangle$ is a group of order exactly $2n$, namely
the dihedral group $D_n$. (Getting an upper bound of $2n$ is easy; the content is that the
group is no smaller.)
\item[(c)] Prove the Nielsen--Schreier theorem: every subgroup of a free group is free. Then prove the
Schreier index formula: if $H \le F$ has finite index $e$ in a free group $F$ of rank $n$,
then $H$ is free of rank $e(n-1) + 1$. Deduce that the free group of rank 2 contains free
subgroups of every finite rank, and of countably infinite rank.
\item[(d)] The word problem asks for an algorithm deciding whether a given word in the generators
represents the identity. Read Novikov's (1955) and Boone's (1958) proof that no such algorithm
exists for all finite presentations, and explain carefully why this does not contradict the fact,
from (a), that every group is a quotient of a free group.
\end{enumerate}

\end{problem}

\noindent Presenting a group by generators and relations is Walther von Dyck's idea
(1882), and it is a bargain with a catch: any group can be written down in a line, and almost
nothing about it can be read off. Nielsen proved the subgroup theorem for finitely generated
subgroups in 1921 by a combinatorial straightening algorithm; Schreier removed the finiteness in
1927; Baer and Levi gave the topological proof in 1936, in which $F$ is the fundamental group of
a bouquet of circles and a subgroup is the fundamental group of a covering graph, so freeness is
just the statement that a graph deformation-retracts to a bouquet. Every known proof uses the
axiom of choice, and provably must. Part (d) is the shock at the end: the objects of (a) are
perfectly explicit, and the questions you want to ask about them are undecidable --- the discovery
that opened combinatorial and then geometric group theory (ALG-25).

\begin{refs}
\item Contexte : Groupe libre --- Wikipédia (\url{https://fr.wikipedia.org/wiki/Groupe_libre})
\item Contexte : Présentation d'un groupe --- Wikipédia (\url{https://fr.wikipedia.org/wiki/Pr%C3%A9sentation_d%27un_groupe})
\item Contexte : Théorème de Nielsen--Schreier --- Wikipédia (\url{https://fr.wikipedia.org/wiki/Th%C3%A9or%C3%A8me_de_Nielsen-Schreier})
\item Lecture : W. Magnus, A. Karrass \& D. Solitar, \textit{Combinatorial Group Theory}, ch. 1--3
\end{refs}

\bigskip

\begin{problem}[{How a group algebra falls apart: Maschke to Wedderburn --- Master}]
\textbf{ALG-40.} \textbf{Problème.} For a field $k$ and a finite group $G$, the group algebra $k[G]$ is the
$k$-vector space with basis $G$ and multiplication extending that of $G$. A module is
\textit{semisimple} if it is a direct sum of simple submodules.
\begin{enumerate}
\item[(a)] Prove Maschke's theorem: if the characteristic of $k$ does not divide $|G|$, then every
$k[G]$-module is semisimple. Then show the hypothesis is needed: for $k = \mathbb{F}_p$ and
$G$ cyclic of order $p$, prove that $k[G]$ is not a direct sum of simple modules.
\item[(b)] Prove Schur's lemma, and deduce that for $k$ algebraically closed of characteristic zero
\[ k[G] \;\cong\; M_{n_1}(k) \times \cdots \times M_{n_t}(k). \]
Deduce that $n_1^2 + \cdots + n_t^2 = |G|$ and that $t$ is the number of conjugacy classes of
$G$ --- and check both against the character table of $S_3$ in ALG-31.
\item[(c)] Prove the Wedderburn--Artin theorem in general: a semisimple ring is isomorphic to a finite
product of matrix rings $M_{n_i}(D_i)$ over division rings, with the $n_i$ and the
isomorphism classes of the $D_i$ uniquely determined.
\item[(d)] Over $\mathbb{R}$ the division rings that can occur are $\mathbb{R}$, $\mathbb{C}$ and
the quaternions $\mathbb{H}$ --- this is Frobenius's 1878 theorem, which you should prove or read.
Then determine the Wedderburn decomposition of $\mathbb{R}[Q_8]$ for the quaternion group
$Q_8$ of ALG-07, and identify which division rings appear.
\end{enumerate}

\end{problem}

\noindent Representation theory begins by refusing to study a group directly: study the
algebra it generates, and let ring theory take it apart. Molien (1893) and Frobenius (1896)
found the decomposition for complex group algebras; Maschke proved the averaging theorem in 1898;
Wedderburn classified finite-dimensional algebras in 1907; and Emil Artin extended the result in
1927 to rings with the descending chain condition, which is why such rings are called Artinian.
Part (a)'s excluded case is not a technicality but a doorway: when $p$ divides $|G|$ the
algebra is no longer semisimple, and Brauer's modular representation theory begins --- the tool
that made the classification of finite simple groups (ALG-22) possible. Part (d) explains why
Hamilton's quaternions cannot be avoided: they are forced by a real group as small as
$Q_8$.

\begin{refs}
\item Contexte : Théorème d'Artin--Wedderburn --- Wikipédia (\url{https://fr.wikipedia.org/wiki/Th%C3%A9or%C3%A8me_d%27Artin-Wedderburn})
\item Contexte : Théorème de Maschke --- Wikipédia (\url{https://fr.wikipedia.org/wiki/Th%C3%A9or%C3%A8me_de_Maschke})
\item Contexte : Théorème de Frobenius (algèbres à division réelles) --- Wikipédia (\url{https://fr.wikipedia.org/wiki/Th%C3%A9or%C3%A8me_de_Frobenius_(alg%C3%A8bre)})
\item Lecture : T. Y. Lam, \textit{A First Course in Noncommutative Rings} (Springer GTM 131), ch. 1--3
\end{refs}

\bigskip

\begin{problem}[{Nakayama's lemma --- Master, short}]
\textbf{ALG-41.} \textbf{Problème.} Let $R$ be a commutative ring with 1, let $J$ be its
Jacobson radical (the intersection of all maximal ideals), and let $M$ be a finitely generated
$R$-module with $JM = M$; prove that $M = 0$. Deduce that if $R$ is local with maximal
ideal $\mathfrak{m}$, and $x_1, \dots, x_n \in M$ have images spanning
$M/\mathfrak{m}M$ as a vector space over $R/\mathfrak{m}$, then $x_1, \dots, x_n$ generate
$M$ --- and show by example that both statements fail without finite generation.
\end{problem}

\noindent Also called the Krull--Azumaya theorem --- Nakayama, who published this form in a
two-page 1951 note, insisted the credit belonged to Krull and Azumaya --- this is the lemma that
makes modules over a local ring behave like vector spaces: a basis of the fibre lifts to a
generating set. Its standard proof is the "determinant trick", a Cayley--Hamilton argument
(LIN-09) applied to a module rather than a space, and its geometric reading is that generators of
a sheaf at a point suffice near the point. Almost every argument in commutative algebra and
algebraic geometry passes through it at some stage.

\begin{refs}
\item Contexte : Lemme de Nakayama --- Wikipédia (\url{https://fr.wikipedia.org/wiki/Lemme_de_Nakayama})
\item Contexte : Radical de Jacobson --- Wikipédia (\url{https://fr.wikipedia.org/wiki/Radical_de_Jacobson})
\item Lecture : M. Atiyah \& I. Macdonald, \textit{Introduction to Commutative Algebra}, ch. 2
\item Lecture : D. Eisenbud, \textit{Commutative Algebra with a View Toward Algebraic Geometry}, ch. 4
\end{refs}

\bigskip


\section*{Recherche}

\begin{problem}[{The inverse Galois problem --- Recherche}]
\textbf{ALG-23.} \textbf{Problème.} Is every finite group the Galois group of some extension of
$\mathbb{Q}$? \textbf{Status: open.} As a first step into the problem:
\begin{enumerate}
\item[(a)] prove that every
finite abelian group is a Galois group over $\mathbb{Q}$ (use cyclotomic fields and
Dirichlet's theorem on primes in arithmetic progressions);
\item[(b)] show that every finite group is
the Galois group of some extension of \textit{some} number field, and explain why this does not
answer the question over $\mathbb{Q}$ itself.
\end{enumerate}

\end{problem}

\noindent Implicit in Hilbert's 1892 work --- where he proved $S_n$ and $A_n$ are
Galois groups over $\mathbb{Q}$ via his irreducibility theorem --- the problem is open after 130
years. Shafarevich proved all solvable groups occur (1954); rigidity methods of Belyi, Matzat,
and Thompson realized many simple groups, including the Monster (Thompson, 1984). A milestone
fell in August 2026: the Mathieu group $M_{23}$, the last of the 26 sporadic groups
outstanding, was realized over $\mathbb{Q}$ by Huang, Jackson, Lee, Poonen, Pries, and Zhang.
The general problem --- even, say, for all finite simple groups --- remains open, and no general
strategy is known.

\begin{refs}
\item Contexte : Théorie de Galois inverse --- Wikipédia (\url{https://fr.wikipedia.org/wiki/Th%C3%A9orie_de_Galois_inverse})
\item Lecture : J.-P. Serre, \textit{Topics in Galois Theory}
\item Article : X. Huang, B. Jackson, K.-H. Lee, B. Poonen, R. Pries, S. Zhang, ``The Mathieu group M$_2$$_3$ is a Galois group over $\mathbb{Q}$'' (2026), arXiv:2608.08538 (\url{https://arxiv.org/abs/2608.08538})
\end{refs}

\bigskip

\begin{problem}[{Burnside problems: bounded exponent --- Recherche}]
\textbf{ALG-24.} \textbf{Problème.} Say a group has \textit{exponent} $n$ if $g^n = 1$ for
every element $g$. Burnside asked (1902): must a finitely generated group of finite exponent
be finite?
\begin{enumerate}
\item[(a)] Prove the case $n = 2$: every group of exponent 2 is abelian, and a finitely
generated one is finite.
\item[(b)] Try $n = 3$ (Burnside settled it in the same 1902 paper).
\item[(c)] Read the status of the rest. \textbf{Status:} the general answer is \textbf{no} --- Novikov and
Adian (1968) constructed infinite finitely generated groups of odd exponent $n \ge 4381$
(Adian later improved the bound to odd $n \ge 665$, and Ivanov (1994) treated all
sufficiently large exponents, even ones included). The \textit{restricted} Burnside problem --- are
there only finitely many \textit{finite} $m$-generator groups of exponent $n$? --- was answered
\textbf{yes} by Zelmanov (1990--91), earning him the 1994 Fields Medal. Sharply open to this day:
is the free Burnside group $B(2,5)$ --- two generators, exponent 5 --- finite or infinite?
Nobody knows.
\end{enumerate}

\end{problem}

\noindent William Burnside posed the question at the very beginning of infinite group
theory, and it shaped a century: the negative solution required the monumental combinatorial
machinery of Novikov--Adian (300+ pages), while Zelmanov's positive solution of the restricted
problem drew on Lie-algebra methods and the odd-order theorem's legacy via the reduction of Hall
and Higman. That $B(2,5)$, an object definable in one line, resists every known technique is a
standing humiliation cheerfully acknowledged by the field.

\begin{refs}
\item Contexte : Problème de Burnside --- Wikipédia (\url{https://fr.wikipedia.org/wiki/Probl%C3%A8me_de_Burnside})
\item Lecture : J. Rotman, \textit{An Introduction to the Theory of Groups}, 4th ed., ch. 12
\item Article : E. Zelmanov, ``Solution of the restricted Burnside problem for groups of odd exponent'' (1990), Izv. Akad. Nauk SSSR Ser. Mat. 54; and ``\dots{}for 2-groups'' (1991), Mat. Sb. 182
\end{refs}

\bigskip

\begin{problem}[{How fast do groups grow? --- Recherche}]
\textbf{ALG-25.} \textbf{Problème.} For a group $G$ generated by a finite set $S$, let
$\gamma(n)$ be the number of elements expressible as products of at most $n$ generators and
inverses --- the \textit{growth function}.
\begin{enumerate}
\item[(a)] Prove that $\mathbb{Z}^d$ has polynomial growth of
degree $d$, and that a free group on two generators has exponential growth.
\item[(b)] Prove that
the growth type (polynomial of given degree, exponential, or between) does not depend on the
choice of $S$.
\item[(c)] Read Gromov's theorem (1981): a finitely generated group has polynomial
growth if and only if it is virtually nilpotent --- and the construction by Grigorchuk (1984) of
a group of \textit{intermediate} growth, strictly between polynomial and exponential.
\textbf{Status of the frontier:} Grigorchuk's Gap Conjecture --- that no group grows strictly
between polynomial growth and $\exp(\sqrt{n})$ --- is \textbf{open}. The precise growth of
Grigorchuk's own group was determined only in 2020, by Erschler and Zheng. It is also still
unknown whether every group of intermediate growth is commensurable with a residually finite
one, and Gromov's theorem keeps being re-proved with new eyes (Kleiner 2010, Ozawa 2018).
\end{enumerate}

\end{problem}

\noindent Growth of groups began with Milnor's questions in the 1960s and Wolf's and
Bass's work on nilpotent groups; Gromov's 1981 solution introduced the audacious idea of viewing a
discrete group from infinitely far away until it becomes a smooth object --- a technique
(Gromov--Hausdorff limits) that reorganized geometry as much as algebra. Grigorchuk's group, built
from automorphisms of an infinite binary tree, showed the landscape between polynomial and
exponential is inhabited; how densely, the Gap Conjecture asks, and forty years on there is no
answer.

\begin{refs}
\item Contexte : Croissance des groupes --- Wikipédia (en anglais) (\url{https://en.wikipedia.org/wiki/Growth_rate_(group_theory)})
\item Contexte : Théorème de Gromov sur les groupes à croissance polynomiale --- Wikipédia (\url{https://fr.wikipedia.org/wiki/Th%C3%A9or%C3%A8me_de_Gromov_sur_les_groupes_%C3%A0_croissance_polynomiale})
\item Contexte : Groupe de Grigorchuk --- Wikipédia (\url{https://fr.wikipedia.org/wiki/Groupe_de_Grigorchuk})
\item Article : A. Erschler, T. Zheng, ``Growth of periodic Grigorchuk groups'' (2020), Invent. Math. 219, 1069--1155, arXiv:1802.09077 (\url{https://arxiv.org/abs/1802.09077})
\end{refs}

\bigskip

\begin{problem}[{State the inverse Galois problem precisely --- Recherche, short}]
\textbf{ALG-32.} \textbf{Problème.} Write a precise statement of the inverse Galois problem --- is every finite group the Galois group of some Galois extension of $ \mathbb{Q} $? --- and summarise its status honestly: open in general, known for all solvable groups by Shafarevich, and known for most finite simple groups, including the monster.
\end{problem}

\noindent Galois theory tells you the group of a given equation; this problem reverses the arrow and asks which groups arise at all, and the reversal turns a computational subject into a wide-open one. That solvable groups and the monster are both settled while the general case resists is a good illustration of how uneven mathematical knowledge can be.

\begin{refs}
\item Contexte : Théorie de Galois inverse --- Wikipédia (\url{https://fr.wikipedia.org/wiki/Th%C3%A9orie_de_Galois_inverse})
\end{refs}

\bigskip

\begin{problem}[{How long is the classification? --- Recherche, short}]
\textbf{ALG-33.} \textbf{Problème.} State the classification of finite simple groups --- every finite simple group is cyclic of prime order, alternating of degree at least five, of Lie type, or one of twenty-six sporadic groups --- and write a short account of why its proof is unusual: it runs to roughly ten thousand journal pages by over a hundred authors, and a "second generation" project to produce a single coherent proof has been under way for decades.
\end{problem}

\noindent This is the largest theorem ever proved, and it raises a real epistemological question the site's readers should confront: what does it mean for a community to know something that no single person has verified end to end? Compare it with the computer-assisted proofs on the Open Problems page --- the difficulty is the same in kind.

\begin{refs}
\item Contexte : Classification des groupes finis simples --- Wikipédia (\url{https://fr.wikipedia.org/wiki/Classification_des_groupes_finis_simples})
\item Voir aussi : Open Problems (\url{open-problems.html})
\end{refs}

\bigskip

\begin{problem}[{Kaplansky's conjectures on group rings --- Recherche}]
\textbf{ALG-42.} \textbf{Problème.} Let $K$ be a field and $G$ a torsion-free group, and form the group ring
$K[G]$ of finite formal sums $\sum k_g\, g$. The \textit{trivial} units are the elements
$k g$ with $k \ne 0$. Kaplansky's three conjectures assert that $K[G]$ has no zero divisors,
no idempotents other than $0$ and $1$, and no units other than the trivial ones.
\textbf{Status:} the unit conjecture --- originally Graham Higman's, popularised by Kaplansky --- is
\textbf{false}, disproved by Giles Gardam in 2021 over $\mathbb{F}_2$; the zero-divisor and
idempotent conjectures remain \textbf{open}, though both are theorems for large classes of groups.
\begin{enumerate}
\item[(a)] Prove that, for fixed $K$ and $G$, absence of zero divisors implies absence of nontrivial
idempotents.
\item[(b)] Prove that torsion must be excluded: if $g \in G$ has finite order $n \ge 2$, exhibit a
zero divisor in $K[G]$, and --- when the characteristic of $K$ does not divide $n$ --- a
nontrivial idempotent.
\item[(c)] Prove that if $G$ carries a total order invariant under left multiplication, then $K[G]$
has no zero divisors and only trivial units. Deduce all three conjectures for free abelian groups
of finite rank and for free groups (ALG-39).
\item[(d)] Study Gardam's counterexample: the group is the fundamental group of the Hantzsche--Wendt flat
3-manifold, a torsion-free crystallographic group, and the field is $\mathbb{F}_2$. Verify by
direct multiplication that his element is a nontrivial unit, and write an account of what the
example does and does not settle --- in particular, why it leaves the zero-divisor conjecture
untouched.
\end{enumerate}

\end{problem}

\noindent Kaplansky's questions, asked in the 1950s and 60s, are the sharpest test of a
simple intuition: a torsion-free group has no finite cycles, so its group ring should have no
"collapse" --- no nilpotents, no idempotents, no exotic units. The intuition held up for eighty
years and resisted every general technique, which made Gardam's 2021 counterexample, found by
reformulating the search as a satisfiability problem and letting a computer hunt through the
finite pieces of a specific crystallographic group, a genuine surprise; a later preprint of his
(2023) reports a counterexample in characteristic zero as well. The remaining two conjectures are
tied to deep machinery elsewhere --- the idempotent conjecture is a cousin of the Baum--Connes
programme in operator algebras --- and their resolution would be major news.

\begin{refs}
\item Contexte : Conjectures de Kaplansky --- Wikipédia (\url{https://fr.wikipedia.org/wiki/Conjecture_de_Kaplansky})
\item Contexte : Anneau de groupe --- Wikipédia (en anglais) (\url{https://en.wikipedia.org/wiki/Group_ring})
\item Article : G. Gardam, ``A counterexample to the unit conjecture for group rings'' (2021), Annals of Mathematics 194, 967--979, arXiv:2102.11818 (\url{https://arxiv.org/abs/2102.11818})
\item Lecture : D. Passman, \textit{The Algebraic Structure of Group Rings}
\end{refs}

\bigskip

\begin{problem}[{The Andrews--Curtis conjecture --- Recherche, short}]
\textbf{ALG-43.} \textbf{Problème.} A presentation $\langle x_1, \dots, x_n \mid r_1, \dots,
r_n \rangle$ with equally many generators and relators is \textit{balanced}; the Andrews--Curtis
conjecture (1965) states that every balanced presentation of the trivial group can be reduced to
$\langle x_1, \dots, x_n \mid x_1, \dots, x_n \rangle$ by finitely many Nielsen transformations
of the relators together with conjugations of relators. Prove it, or produce a counterexample ---
\textbf{status: open}, with no counterexample known and a widespread expectation that the
conjecture is false.
\end{problem}

\noindent Andrews and Curtis raised the question in a 1965 paper on free groups and
handlebodies, and it is a rare open problem whose statement needs nothing beyond the presentations
of ALG-39: a purely combinatorial game on words that has resisted sixty years of attack. It has
topological teeth --- the Zeeman collapsibility conjecture implies it, and it stands in the way of
certain approaches to smooth four-dimensional problems --- and it is a favourite testing ground for
automated search, since celebrated candidate counterexamples have one after another been
trivialised by computer once enough moves were allowed.

\begin{refs}
\item Contexte : Conjecture d'Andrews--Curtis --- Wikipédia (\url{https://fr.wikipedia.org/wiki/Conjecture_d%27Andrews-Curtis})
\item Article : J. J. Andrews \& M. L. Curtis, ``Free groups and handlebodies'' (1965), Proceedings of the American Mathematical Society 16, 192--195
\end{refs}

\bigskip


\section*{Pour aller plus loin}


\vfill
\noindent\emph{Hors de la Caverne} --- des probl\`emes, pas de r\'eponses.
\end{document}
